\subsection{The exact harmonic Jacobian and quadratic norm comparison}

\begin{statement}{A.108}{Proved}
{The harmonic joint character factors through an elliptic norm cocycle}
Fix $R\in\Q\setminus\{0,1,\tfrac12\}$, and put
$F(T)=T(T-1)(T-\tfrac12)$.
Let $C_4,C_2,D$ be the smooth projective normalizations of
\[
\begin{split}
 C_4:\quad&T-c_i=(R-c_i)X_i^4,\\
 C_2:\quad&T-c_i=(R-c_i)U_i^2
                 \qquad(c_i=0,1,\tfrac12),\\
 D:\quad&W^4=F(T)/F(R).
\end{split}
\tag{A.108.1}
\]
All three have their indicated rational point over $T=R$
with every root coordinate one.
The actual maps are
$\pi:C_4\to D$, $W=\prod_iX_i$, and
$r:C_4\to C_2$, $U_i=X_i^2$.

Choose $d,y_0\in\Q^\times$ with
$y_0^2=d^3u(u^2-1)$, where $u=2R-1$.
Such a choice always exists: take
$d=u(u^2-1)$ and $y_0=d^2$.
Put
\[
 E:\ y^2=x^3-d^2x,\qquad P_R=(du,y_0),\qquad
 q:D\longrightarrow E,\quad
       (x,y)=\bigl(d(2T-1),y_0W^2\bigr).
\tag{A.108.2}
\]
There is a unique pointed $\Q$-isomorphism
$\psi:C_2\to E$, taking the old point to $O$, for which
the map
$(T,U_i)\mapsto(d(2T-1),y_0\prod_iU_i)$ is
$[2]\psi+P_R$.
Write $\bar r=\psi r$ and $J_4=\operatorname{Jac}(C_4)$.

For $Q\in E(\Q)$ define the continuous cocycle
\[
 \kappa_Q(g)=
 \begin{cases}
 0,&g(i)=i,\\
 [i]Q,&g(i)=-i,
 \end{cases}
 \qquad [i](x,y)=(-x,iy).
\tag{A.108.3}
\]
The normalized Brauer group of $C_4$ is identified with
$H^1(\Q,J_4)$ by its rational base point and the
Hochschild--Serre map of A.88.
Under this exact identification the homomorphism
\[
 \Theta_R:E(\Q)\longrightarrow
       \ker\{\operatorname{Br}(C_4)\to\operatorname{Br}(\Q)\},
 \qquad
 \epsilon(\Theta_R(Q))=\bar r^{\,*}[\kappa_Q],
\tag{A.108.4}
\]
has the following explicit symbols.
For
\[
 Q=(-dv,\rho y_0),\qquad
 c=(1+v)/2\notin\{0,1,\tfrac12,R\},
                     \qquad \rho\ne0,
\tag{A.108.5}
\]
one has
\[
 \Theta_R(Q)=
        \left(\frac{T-c}{R-c},\,2\rho\right)
                            \quad\hbox{in }\operatorname{Br}(C_4).
\tag{A.108.6}
\]
In particular this class restricts to zero over $\Q(i)$.

The homomorphism $\Theta_R$ kills $E[2](\Q)$, $P_R$,
and every rational point whose nonzero $x$-coordinate
is a positive or negative rational square.
More precisely it factors through the actual squareclass map
\[
 \begin{gathered}
 \frac{\delta_x(E(\Q))}
       {\langle-1,d,x(P_R)\rangle}
 \longrightarrow \operatorname{Br}(C_4),\\
 \delta_x(O)=1,\qquad
 \delta_x((0,0))=-1,\\
 \delta_x(Q)=x(Q)\bmod\Q^{\times2}
                         \quad(x(Q)\ne0).
 \end{gathered}
\tag{A.108.7}
\]
All denominator classes in (A.108.7) belong to
$\delta_x(E(\Q))$.

For any characteristic-zero field $k$, the source class
$[\kappa_Q]\in H^1(k,E)$ is zero exactly when
$Q$ is a norm from $E(k(i))$; if $i\in k$, use the split
quadratic algebra and its sum map.
For $Q=(a,b)\ne O$ this norm equation has the explicit
projective genus-one model with affine equation
\[
 z^2=-m^4+6am^2+8bm+3a^2-4d^2.
\tag{A.108.8}
\]
The conclusion about the Brauer class is the actual
pushforward (A.108.4); no injectivity of
$H^1(k,E)\to H^1(k,J_4)$ is asserted.
\end{statement}
\begin{proof}
\emph{The finite maps and their Jacobian kernels.}
Over $\overline\Q(T)$ the three classes $T-c_i$
are independent modulo fourth powers, by their
valuations at the three distinct $c_i$.
Thus $C_4$ has degree $64$ over the $T$-line, and
$D$ has degree four.
The homomorphism on their deck groups is
\[
       \mu_4^3\longrightarrow\mu_4,\qquad
                    (\zeta_0,\zeta_1,\zeta_{1/2})
                         \longmapsto\prod_i\zeta_i.
\tag{A.108.9}
\]
At each of the four boundary points the cyclic
inertia map is injective: its product exponent is
one at a finite branch point and $-3$ modulo four
at infinity.  Hence $\pi$ extends to an etale cover
of the proper normalizations, with group $\mu_4^2$.
Riemann--Hurwitz gives $g(D)=3$ and $g(C_4)=33$.

Write $B_0,B_1,B_{1/2},B_\infty$ for the four
rational branch points of $D$.
For $L_i=[B_i-B_\infty]\in J_D=\operatorname{Jac}(D)$,
\[
 4L_i=0,\qquad L_0+L_1+L_{1/2}=0,\qquad
 \ker(\pi^*:J_D\to J_4)
       =\langle L_0,L_1,L_{1/2}\rangle
                          \simeq(\Z/4)^2.
\tag{A.108.10}
\]
Here the divisors of $T-c_i$ and $W$ give the first
two relations literally.
For the last equality, a line bundle trivialized
by an etale connected Galois cover acquires, on
its trivialization, a deck action by a character:
the units on the proper geometrically connected
cover are constants.
Conversely such a character descends the trivial
line bundle.  Thus the pullback kernel is exactly
the character group of $\mu_4^2$, of order sixteen.
The three fourth-root extensions in (A.108.1)
are precisely the character line bundles $L_i$;
their one relation is the displayed product relation.
This proves (A.108.10), including its order.

Over $k_0=\Q(i)$ let
$a_D:(T,W)\mapsto(T,iW)$, and denote its induced
endomorphism of $J_D$ by $a$.
The three regular differentials
$dT/W^2,dT/W^3,TdT/W^3$ form a basis.
Their $a_D$-eigenvalues are $-1,i,i$.
On integral first homology over $\C$ the eigenvalues
are therefore $-1,-1,i,i,-i,-i$.
The endomorphism $1-a$ has nonzero determinant
\[
             2^2(1-i)^2(1+i)^2=16.
\tag{A.108.11}
\]
For a complex torus the degree of an isogeny is
the index of its image on the integral lattice,
so (A.108.11) gives $\deg(1-a)=16$.
This computation is unchanged by the characteristic-zero
base-change comparison
\cite[Remark 3.3(c), Theorems 21.3 and 21.5]{AgentMilneLEC}.
Every $L_i$ is fixed by $a$.
Equation (A.108.10) and equality of orders consequently give
\[
                    \ker(1-a)=\ker\pi^*.
\tag{A.108.12}
\]
The quotient by this kernel supplies an actual
closed immersion $j:J_D\to J_4$ over $k_0$ such that
\[
                  \pi^*=j(1-a).
\tag{A.108.13}
\]
This is a factorization of homomorphisms of
abelian varieties, not a division of cohomology classes.

\emph{The rational elliptic factor and the square-coordinate map.}
Let a bar denote conjugation of $k_0/\Q$.
Since $\bar a=a^{-1}$ and $\pi^*$ is rational,
cancelling the isogeny $1-a^{-1}$ in (A.108.13) gives
\[
              \bar j=-ja.
\tag{A.108.14}
\]
The map $q$ in (A.108.2) is the quotient by
$a_D^2$, and $a_D$ induces the negation map on $E$.
Thus $a q^*=-q^*$, and the homomorphism
$h=j q^*:E\to J_4$ is defined over $\Q$ by
(A.108.14).

The cover $C_2\to\mathbf P^1_T$ has group $\mu_2^3$,
degree eight and four inertia groups of order two.
Its genus is one.
The product-of-roots map $C_2\to E$ has degree four
and is etale, because its inertia maps are injective
at all four boundary points.
Its deck transformations are the four sign vectors
of product one.
With the old rational point as origin, each
nonidentity deck involution is a translation:
an automorphism of an elliptic curve with nontrivial
linear part has a geometric fixed point, since
one minus that linear part is an isogeny.
These three fixed-point-free involutions and the
identity are therefore exactly the full rational
two-torsion translation subgroup.
Subtracting $P_R$ from the target map produces a
pointed isogeny with kernel $C_2[2]$.
It factors uniquely through $[2]$, giving the
$\Q$-isomorphism $\psi$ in the statement.

Translation has trivial pullback on degree-zero
Picard groups.  Consequently the two maps from
$C_4$ to $E$ give the literal homomorphism identity
\[
      2\bar r^{\,*}=\pi^*q^*
                =j(1-a)q^*=2j q^*=2h.
\tag{A.108.15}
\]
The group of homomorphisms of abelian varieties
has no additive torsion: a nonzero homomorphism
has a positive-dimensional image and cannot be
killed by a nonzero integer.
Cancelling two proves
\[
                         h=\bar r^{\,*}.
\tag{A.108.16}
\]
This identifies the specific elliptic pullback
appearing in (A.108.4).

\emph{The divisor cocycle and its descent over $\Q$.}
For (A.108.5), the equation of $E$ gives
$v(v^2-1)=-\rho^2u(u^2-1)$.
Thus $W^4=-\rho^2$ at $T=c$.
Put $s=\sqrt{2\rho}$, and let $Z_+$ be the two
points in that fibre satisfying
$W^2-sW+\rho=0$.
For $D_c=Z_+-2B_\infty$ and
$f=(T-c)/(R-c)$ one has
\[
              D_c+\sigma D_c=\operatorname{div}(f),
\tag{A.108.17}
\]
where $\sigma(s)=-s$.
The quaternion $(f,2\rho)$ is unramified on $D$:
at $T=c$ the element $(1-W^2/\rho)W$ squares to
$2\rho$, and its only pole is of order four at
$B_\infty$.
Its value at the old point is zero.
As in A.88, the cocycle $\eta(g)$ equal to $0$
or $[D_c]$ according as $g$ fixes $s$ or not
has Hochschild--Serre image $-[\eta]$ for this
quaternion, using the same total-complex sign.
The rational-point splitting and this divisor
description are
\cite[Propositions 4.3.2 and 4.3.4, Remarks 4.3.3,
4.3.5 and 4.5.2]{AgentCTSkoro2019}.

The following computation works even when
$\Q(s)$ equals $\Q$ or $\Q(i)$.
On $D_{\Q(i,s)}$ choose
\[
 P=(c,s(1+i)/2),\qquad z=[P-B_\infty]\in J_D,
                    \qquad q(P)=[i]Q.
\tag{A.108.18}
\]
The other point of $Z_+$ is $a_D^{-1}P$.
Hence $[D_c]=(1+a^{-1})z$.
For $g\in G_\Q$ write
$g(i)=(-1)^\epsilon i$ and $g(s)=(-1)^\delta s$.
Then
\[
 gz=a^{2\delta-\epsilon}z,\qquad
 g(a)=a^{(-1)^\epsilon},\qquad
 g(j)=j(-a)^\epsilon.
\tag{A.108.19}
\]
Take the zero-cochain $b=jaz\in J_4(\overline\Q)$.
For $\epsilon=0$,
$\pi^*\eta(g)=gb-b$.
For $\epsilon=1$, both choices of $\delta$ give
\[
 \begin{split}
 \pi^*\eta(g)-(gb-b)
     &=j(a^{-1}+a)z\\
     &=j a^{-1}(1+a^2)z
       =-h([i]Q).
 \end{split}
\tag{A.108.20}
\]
The last equality uses the actual quotient divisor
identity
$(1+a^2)z=q^*([i]Q-O)$ and the action $a^{-1}=-1$
on $q^*E$.
Thus $-[\pi^*\eta]=h[\kappa_Q]$ with the sign
convention above.
Pullback of Brauer classes commutes with the
Hochschild--Serre map.
Both classes are normalized at the old point,
where $f=1$; the split comparison is injective
on this normalized group.
Equations (A.108.16)--(A.108.20) prove (A.108.6)
as an equality in $\operatorname{Br}(C_4)$.
They also prove its vanishing after restriction
to $\Q(i)$.

\emph{The exact killed directions.}
The map $Q\mapsto\kappa_Q$ is additive.
It is a cocycle because conjugation sends
$[i]Q$ to $-[i]Q$.
Twice it is the coboundary of $-[i]Q$.
These facts prove that $\Theta_R$ is a
homomorphism with two-torsion image.

For any nonidentity rational two-torsion point
$T_i=q(B_i)$, ramification of the degree-two map
$q$ gives
\[
              q^*(T_i-O)=2L_i,\qquad
    \bar{(-jL_i)}-(-jL_i)=2jL_i=hT_i.
\tag{A.108.21}
\]
Here $aL_i=L_i$ and (A.108.14) were used.
For the cocycle of a specified $T_i$, choose the
index $k$ with $q(B_k)=[i]T_i$.
Then $-jL_k$ is its explicit coboundary primitive.
This proves $\Theta_R(E[2])=0$.

Now let $Q=(a,b)$ have $a\ne0$ and
$a=z_0^2$ in $\Q(i)$; for rational $a$ this says
$a$ is a positive or negative rational square.
The explicit point
\[
 w=i(a+b/z_0),\qquad
             S=(w,(1+i)wz_0)\in E(\Q(i))
\tag{A.108.22}
\]
satisfies $[1+i]S=Q$.
Indeed $w^2-2iaw-d^2=0$ proves the curve equation,
and the chord formula for addition with $[i]S$
gives the two coordinates of $Q$.
Conjugating $[1+i]S=Q$ shows
\[
             \bar S-iS\in\ker[1-i]=\{O,(0,0)\}.
\tag{A.108.23}
\]
Thus $\bar S-S=[i]Q+\varepsilon(0,0)$.
After applying $h$, the last two-torsion term is
a coboundary by (A.108.21).
This proves $\Theta_R(Q)=0$, including the possible
two-torsion correction in (A.108.23).

For the source $Q=P_R$ in (A.108.5), one has
$c=1-R$ and $\rho=1$.
The third old equation gives
\[
 \frac{T-c}{R-c}=\frac{X_{1/2}^4+1}{2},\qquad
 X_{1/2}^4+1=
 N_{\Q(\sqrt2)/\Q}
        (X_{1/2}^2+1+\sqrt2X_{1/2}),\qquad
                  N(2+\sqrt2)=2.
\tag{A.108.24}
\]
Hence this quaternion is split and
$\Theta_R(P_R)=0$.

The map $\delta_x$ in (A.108.7) is a homomorphism.
For a line $y=mx+b_0$ intersecting $E$ in three
points, the product of their $x$-coordinates is
$b_0^2$, by the constant coefficient of
$x^3-d^2x-(mx+b_0)^2$.
This proves the addition formula modulo squares,
including tangencies by the same polynomial.
The local parameter at $(0,0)$ gives its value
$-d^2=-1$ modulo squares; the identity has value one.
The other two two-torsion points have values
$d$ and $-d$.
Combining this homomorphism with the killed
square, two-torsion and base-point directions
proves precisely the factorization (A.108.7).

\emph{The source norm torsor and its actual quartic equation.}
Suppose first that $i\notin k$ and put $K=k(i)$.
If $\kappa_Q$ is a coboundary, a point trivializing
it is fixed by $G_K$, hence lies in $E(K)$.
Writing $\bar S-S=[i]Q$ gives
\[
          N_{K/k}([i]S)=[i](S-\bar S)=Q.
\tag{A.108.25}
\]
Conversely a point $A\in E(K)$ of norm $Q$
makes $-[i]A$ a coboundary primitive.
This proves the source norm criterion.
For split $K$ the norm is the sum map and is
onto, agreeing with the zero cocycle.

For a nonidentity source $Q=(a,b)$, a finite
line through $-Q$ has equation
$y=mx-b-ma$.
Factoring its intersection polynomial by $x-a$
leaves
\[
 x^2+(a-m^2)x+(a^2+am^2-d^2+2mb),
\tag{A.108.26}
\]
whose discriminant is
$m^4-6am^2-8bm-3a^2+4d^2$.
The two remaining intersections are conjugate
over $k(i)$ exactly when this discriminant is
$-z^2$; their sum is $Q$ by the chord law.
The case $z=0$ gives a rational half of $Q$ and
is included.
Conversely a pair $A,\bar A$ of sum $Q\ne O$
has this line, with the tangent interpretation
when the points coincide.
The two points at infinity on the quartic
require a square root of $-1$.
Thus its projective model gives exactly
(A.108.8) and the stated norm criterion.
\end{proof}

The comparison (A.108.4) retains the possible kernel of
$H^1(k,E)\to H^1(k,J_4)$.
The elliptic Jacobian homomorphism itself is the
specific pullback $\bar r^{\,*}$ of (A.108.16).
Consequently proving a source norm obstruction does
not by itself prove a nonconstant quaternion on $C_4$.
The Jacobian homomorphism $\bar r^{\,*}$ itself has
zero kernel: the degree-eight cover $r:C_4\to C_2$
has group $\mu_2^3$, and its inertia groups over the
three finite old boundary fibres are the three
coordinate sign changes, which generate that group.
Any line bundle in the pullback kernel would give
a character of this group trivial on every inertia
group, by the same descent argument as (A.108.10).
There is no such nontrivial character.
This injectivity on abelian varieties does not
remove the possible kernel on $H^1$: in the exact
sequence with quotient $J_4/\bar r^{\,*}E$, rational
points of that quotient have a connecting map into
$H^1(k,E)$.
The exact independent-squareclass and local searches
are retained in
\path{code/harmonic_descent_point_search.py} and
\path{code/harmonic_independent_descent_search.py},
with their outputs.
The separate exact certificate
\path{code/harmonic_jacobian_descent_certificate.py},
with retained output, checks the uniform Galois signs,
the finite branch kernel, both coordinates of the
$[1+i]$ division formula, and the norm identities
and quartic factorization used in the proof.
